% Options for packages loaded elsewhere
\PassOptionsToPackage{unicode}{hyperref}
\PassOptionsToPackage{hyphens}{url}
\PassOptionsToPackage{dvipsnames,svgnames,x11names}{xcolor}
\documentclass[
  10pt,
  fontset=fandol]{ctexart}
\usepackage{xcolor}
\usepackage[margin=16mm]{geometry}
\usepackage{amsmath,amssymb}
\setcounter{secnumdepth}{-\maxdimen} % remove section numbering
\usepackage{iftex}
\ifPDFTeX
  \usepackage[T1]{fontenc}
  \usepackage[utf8]{inputenc}
  \usepackage{textcomp} % provide euro and other symbols
\else % if luatex or xetex
  \usepackage{unicode-math} % this also loads fontspec
  \defaultfontfeatures{Scale=MatchLowercase}
  \defaultfontfeatures[\rmfamily]{Ligatures=TeX,Scale=1}
\fi
\usepackage{lmodern}
\ifPDFTeX\else
  % xetex/luatex font selection
\fi
% Use upquote if available, for straight quotes in verbatim environments
\IfFileExists{upquote.sty}{\usepackage{upquote}}{}
\IfFileExists{microtype.sty}{% use microtype if available
  \usepackage[]{microtype}
  \UseMicrotypeSet[protrusion]{basicmath} % disable protrusion for tt fonts
}{}
\makeatletter
\@ifundefined{KOMAClassName}{% if non-KOMA class
  \IfFileExists{parskip.sty}{%
    \usepackage{parskip}
  }{% else
    \setlength{\parindent}{0pt}
    \setlength{\parskip}{6pt plus 2pt minus 1pt}}
}{% if KOMA class
  \KOMAoptions{parskip=half}}
\makeatother
\usepackage{longtable,booktabs,array}
\newcounter{none} % for unnumbered tables
\usepackage{calc} % for calculating minipage widths
% Correct order of tables after \paragraph or \subparagraph
\usepackage{etoolbox}
\makeatletter
\patchcmd\longtable{\par}{\if@noskipsec\mbox{}\fi\par}{}{}
\makeatother
% Allow footnotes in longtable head/foot
\IfFileExists{footnotehyper.sty}{\usepackage{footnotehyper}}{\usepackage{footnote}}
\makesavenoteenv{longtable}
\setlength{\emergencystretch}{3em} % prevent overfull lines
\providecommand{\tightlist}{%
  \setlength{\itemsep}{0pt}\setlength{\parskip}{0pt}}
\usepackage{amsmath,amssymb,mathtools,needspace}
\xeCJKDeclareCharClass{CJK}{"2032,"0400->"04FF,"2160->"216B,"2460->"2473,"25A0,"25C7,"25CB,"25CF}
\IfFontExistsTF{Arial Unicode MS}{\xeCJKsetup{AutoFallBack=true}\setCJKfallbackfamilyfont{\CJKrmdefault}{Arial Unicode MS}}{}
\setcounter{secnumdepth}{0}
\setlength{\emergencystretch}{2em}
\allowdisplaybreaks
\usepackage{bookmark}
\IfFileExists{xurl.sty}{\usepackage{xurl}}{} % add URL line breaks if available
\urlstyle{same}
\hypersetup{
  pdftitle={插值型的求导公式},
  colorlinks=true,
  linkcolor={Maroon},
  filecolor={Maroon},
  citecolor={Blue},
  urlcolor={Blue},
  pdfcreator={LaTeX via pandoc}}

\title{插值型的求导公式}
\author{}
\date{}

\begin{document}
\maketitle

\subsubsection{4.5.2 插值型的求导公式}\label{ux63d2ux503cux578bux7684ux6c42ux5bfcux516cux5f0f}

对于列表函数 \(y=f(x)\)（见表 4.7），运用插值原理，可以建立插值多项式 \(y=P_n(x)\) 作为它的近似。由于多项式的求导比较容易，取 \(P_n'(x)\) 的值作为 \(f'(x)\) 的近似值，这样建立的数值公式

\[
f'(x)\approx P_n'(x)
\tag{4.5.1}
\]

统称为插值型的求导公式。

\textbf{表 4.7}

{\def\LTcaptype{none} % do not increment counter
\begin{longtable}[]{@{}cccccc@{}}
\toprule\noalign{}
\(x\) & \(x_0\) & \(x_1\) & \(x_2\) & \(\cdots\) & \(x_n\) \\
\midrule\noalign{}
\endhead
\bottomrule\noalign{}
\endlastfoot
\(y\) & \(y_0\) & \(y_1\) & \(y_2\) & \(\cdots\) & \(y_n\) \\
\end{longtable}
}

必须指出，即使 \(f(x)\) 与 \(P_n(x)\) 的值相差不多，导数的近似值 \(P_n'(x)\) 与导数的真值 \(f'(x)\) 仍然可能差别很大，因而在使用求导公式（4.5.1）时应特别注意误差的分析。

依据插值余项定理，求导公式（4.5.1）的余项为

\[
f'(x)-P_n'(x)=\frac{f^{(n+1)}(\xi)}{(n+1)!}\omega_{n+1}'(x)
+\frac{\omega_{n+1}(x)}{(n+1)!}\frac{\mathrm{d}}{\mathrm{d}x}f^{(n+1)}(\xi),
\]

其中

\[
\omega_{n+1}(x)=\prod_{k=0}^{n}(x-x_i).
\]

在此余项公式中，由于 \(\xi\) 是 \(x\) 的未知函数，无法对其第二项 \(\dfrac{\omega_{n+1}(x)}{(n+1)!}\dfrac{\mathrm{d}}{\mathrm{d}x}f^{(n+1)}(\xi)\) 作出进一步的说明，因此，对于随意给出的点 \(x\)，误差 \(f'(x)-P_n'(x)\) 是无法预估的。但是，如果限定求某个节点 \(x_k\) 上的导数值，那么上面的第二项因 \(\omega_{n+1}(x_k)=0\) 而变为

（本句续下页。）

\begin{quote}
校注（agent 补充，CH04B-E004）：原页``其中''后的节点多项式印为 \(\prod_{k=0}^{n}(x-x_i)\)，乘积指标 \(k\) 与因子下标 \(i\) 不一致，已按原扫描保留。结合本页表 4.7 和下文 \(\omega_{n+1}(x_k)=0\)，应将因子写作 \((x-x_k)\)，或统一将乘积指标改为 \(i\)；这是原书疑似下标排印错误，非 OCR 错误。\href{../assets/ch04b/review-p113-product.png}{局部原扫描}。
\end{quote}

\href{../../source-images/pdf-114.jpeg}{原页图像}

\subsubsection{4.5.2 插值型的求导公式（续）}\label{ux63d2ux503cux578bux7684ux6c42ux5bfcux516cux5f0fux7eed}

零，这时有余项公式

\[
f'(x_k)-P_n'(x_k)=\frac{f^{(n+1)}(\xi)}{(n+1)!}\omega_{n+1}'(x_k).
\tag{4.5.2}
\]

下面仅仅考察节点处的导数值。为简化讨论，假定所给的节点是等距的。

\paragraph{1. 两点公式}\label{ux4e24ux70b9ux516cux5f0f}

设已给出两个节点 \(x_0,x_1\) 上面的函数值 \(f(x_0),f(x_1)\)，作线性插值公式

\[
P_1(x)=\frac{x-x_1}{x_0-x_1}f(x_0)+\frac{x-x_0}{x_1-x_0}f(x_1).
\]

对上式两端求导，记 \(x_1-x_0=h\)，有

\[
P_1'(x)=\frac1h[-f(x_0)+f(x_1)],
\]

于是有下列求导公式：

\[
P_1'(x_0)=\frac1h[f(x_1)-f(x_0)],\quad
P_1'(x_1)=\frac1h[f(x_1)-f(x_0)].
\]

而利用余项公式（4.5.2）知，带余项的两点公式是

\[
f'(x_0)=\frac1h[f(x_1)-f(x_0)]-\frac h2f''(\xi),
\]

\[
f'(x_1)=\frac1h[f(x_1)-f(x_0)]+\frac h2f''(\xi).
\]

\paragraph{2. 三点公式}\label{ux4e09ux70b9ux516cux5f0f}

设已给出三个节点 \(x_0,x_1=x_0+h,x_2=x_0+2h\) 上的函数值，作二次插值

\[
\begin{aligned}
P_2(x)={}&\frac{(x-x_1)(x-x_2)}{(x_0-x_1)(x_0-x_2)}f(x_0)
+\frac{(x-x_0)(x-x_2)}{(x_1-x_0)(x_1-x_2)}f(x_1)\\
&+\frac{(x-x_0)(x-x_1)}{(x_2-x_0)(x_2-x_1)}f(x_2),
\end{aligned}
\]

令 \(x=x_0+th\)，上式可表示为

\[
P_2(x_0+th)=\frac12(t-1)(t-2)f(x_0)-t(t-2)f(x_1)
+\frac12t(t-1)f(x_2),
\]

两端对 \(t\) 求导，有

\[
P_2'(x_0+th)=\frac1{2h}[(2t-3)f(x_0)-(4t-4)f(x_1)+(2t-1)f(x_2)].
\tag{4.5.3}
\]

这里撇号表示对变量 \(x\) 求导数。上式分别取 \(t=0,1,2\)，得到以下三种三点公式：

\[
P_2'(x_0)=\frac1{2h}[-3f(x_0)+4f(x_1)-f(x_2)],
\]

\[
P_2'(x_1)=\frac1{2h}[-f(x_0)+f(x_2)],
\]

（第三个公式续下页。）

\href{../../source-images/pdf-115.jpeg}{原页图像}

\subsubsection{4.5.2 插值型的求导公式（续）}\label{ux63d2ux503cux578bux7684ux6c42ux5bfcux516cux5f0fux7eed-1}

\paragraph{2. 三点公式（续）}\label{ux4e09ux70b9ux516cux5f0fux7eed}

\[
P_2'(x_2)=\frac1{2h}[f(x_0)-4f(x_1)+3f(x_2)],
\]

而带余项的三点求导公式为

\[
f'(x_0)=\frac1{2h}[-3f(x_0)+4f(x_1)-f(x_2)]+\frac{h^2}3f'''(\xi),
\]

\[
f'(x_1)=\frac1{2h}[-f(x_0)+f(x_2)]-\frac{h^2}6f'''(\xi),
\tag{4.5.4}
\]

\[
f'(x_2)=\frac1{2h}[f(x_0)-4f(x_1)+3f(x_2)]+\frac{h^2}3f'''(\xi),
\]

其中式（4.5.4）是我们所熟悉的中点公式。在三点公式中，它由于少用了一个函数值 \(f(x_1)\) 而引人注目。

用插值多项式 \(P_n(x)\) 作为 \(f(x)\) 的近似函数，还可以建立高阶数值微分公式

\[
f^{(k)}(x)\approx P_n^{(k)}(x),\quad k=1,2,\cdots.
\]

例如，将式（4.5.3）再对 \(t\) 求导一次，有

\[
P_2''(x_0+th)=\frac1{h^2}[f(x_0)-2f(x_1)+f(x_2)],
\]

于是有

\[
P_2''(x_1)=\frac1{h^2}[f(x_1-h)-2f(x_1)+f(x_1+h)],
\]

而带余项的二阶三点公式为

\[
f''(x_1)=\frac1{h^2}[f(x_1-h)-2f(x_1)+f(x_1+h)]-\frac{h^2}{12}f^{(4)}(\xi).
\tag{4.5.5}
\]

\subsection{本节覆盖原页的校注记录}\label{ux672cux8282ux8986ux76d6ux539fux9875ux7684ux6821ux6ce8ux8bb0ux5f55}

下列记录按原页关联，含同页相邻小节的校注；计算时同时读取上文校注和完整原页。

\begin{itemize}
\tightlist
\item
  \texttt{CH04B-E004}：原书 PDF 页 113
\end{itemize}

\end{document}
